%%%%%%%%%%%%%%%%%%%%%%%%%%%%%%%%%%%%%%%%%%%%%%%%%%%%%%%%%%%%%%%%%%%%%%%%%%%%%%%
%% LaTeX-Vorlage für Abschlussarbeiten                                       %%
%% (TH Köln -Campus Gummersbach, Fak. 10)                                    %%
%%                                                                           %%
%% Gemäß dem Merkblatt zur Anfertigung von Projekt-, Bachelor-, Master- und  %%
%% Diplomarbeiten der Fakultät 10 von Frau Prof. Dr. Halfmann &              %%
%% Herr Prof. Dr. Rühmann (Version vom 27.01.2008)                           %%
%%                                                                           %%                                                                            
%% Bitte sprechen Sie unbedingt mit Ihrer Betreuerin bzw. Ihrem Betreuer     %%
%% bezüglich der Ausgestaltung Ihrer Arbeit!                                 %%
%%                                                                           %%
%%                                                                           %%
%% MERKKASTEN IN DIESER VORLAGE:                                             %%
%% In dieser Vorlage finden Sie Merkkasten, die Ihnen Informationen          %%
%% zu bestimmten, formalen Aspekten geben. Sprechen Sie immer auch mit       %% 
%% Ihrer Betreuerin bzw. Ihrem Betreuer dazu an.                             %%                       
%% Für die eigene Verwendung der Vorlage entfernen oder kommentieren Sie die %%
%% Merkkasten. Die betreffenden Bereiche für die Merkkasten in der Vorlage   %%
%% sind wie folgt kommentiert: <MERKKASTEN> ... </MERKKASTEN>.               %%                            %%                                                                           %%
%%                                                                           %%
%% LIZENZ:                                                                   %%
%% Diese Vorlage darf nicht kommerziell verbreitet                           %%
%% werden. Eine nicht-kommerzielle Weitergabe ist                            %% 
%% gestattet.                                                                %%
%%                                                                           %%
%% Von Ludger Schönfeld, M. Sc.,
%% 2014-2017                            %%
%%%%%%%%%%%%%%%%%%%%%%%%%%%%%%%%%%%%%%%%%%%%%%%%%%%%%%%%%%%%%%%%%%%%%%%%%%%%%%%

%%%%%%%%%%%%%%%%%%%%%%%%%%%%%%%%%%%%%%%%%%%%%
%% HEADER                                  %%
%%%%%%%%%%%%%%%%%%%%%%%%%%%%%%%%%%%%%%%%%%%%%
\documentclass[a4paper,12pt,oneside]{article}
% Optionen:
% - a4paper => DIN A4-Format
% - 12pt    => Schriftgröße (weitere  
%              grundlegende Fontgrößen: 10pt, 11pt)
% - oneside => Einseitiger Druck

%% Verwendete Pakete:
\usepackage[ngerman]{babel} % für die deutsche Sprache
\usepackage{caption} % Für schönere Bildunterschriften
\usepackage[T1]{fontenc} % Schriftkodierung (Für Sonderzeichen u.a.)
\usepackage[utf8]{inputenc} % Für die direkte Eingabe von Umlauten im Editor u.a.
\usepackage{fancyhdr} % Für Kopf- und Fußzeilen
\usepackage{lscape} % Für Querformat

%% Schriften (Beispiele)
%% Weitere LaTeX-Schriften im "LaTeX Font Catalogue"
%% unter: http://www.tug.dk/FontCatalogue/.
%% ACHTUNG: Ggf. müssen Schriften noch installiert 
%% werden!

% Serifen-Schriften:
\usepackage{lmodern} % Schriftart "Latin Modern"
%\usepackage{garamond} % Schriftart "Garamond"

%Sans Serif-Schriften:
%\usepackage[scaled]{uarial}
%\usepackage[scaled]{helvet}
%%--------------
\usepackage[normalem]{ulem} % Für das Unterstreichen von Text z.B. mit \uline{}
\usepackage[left=3cm,right=2cm,top=1.5cm,bottom=1cm,
textheight=245mm,textwidth=160mm,includeheadfoot,headsep=1cm,
footskip=1cm,headheight=14.599pt]{geometry} % Einrichtung der Seite 

\usepackage{graphicx} % Zum Laden von Graphiken
% INFO: Graphiken einbinden
%
% \includegraphics[scale=1.00]{dateiname}
%
% => Ausgabeformat: PDF-Dokument:
%    Es können die folgenden (Graphik-)formate eingebunden
%    werden: .jpg, .png, .pdf, .mps
% 
% => Ausgabeformat: DVI/PS:
%    Folgende (Graphik-)formate werden unterstützt:
%    .eps, .ps, .bmp, .pict, .pntg
\usepackage{epstopdf}

% Pakete für Tabellen
\usepackage{tabularx} % Einfache Tabellen
\usepackage{longtable} % Tabellen als Gleitobjekte (für die Aufteilung bei langen 
 %Tabellen über mehrere Seiten)
\usepackage{multirow} % Für das Verbinden von Zeilen innerhalb einer Tabelle mit
 % \multirow{anzahl}{*}{Text}

% (Zusatz-)Pakete für Formeln
\usepackage{amsmath}
\usepackage{amsthm}
\usepackage{amsfonts}

\usepackage{setspace} % Paket zum Setzen des Zeilenabstandes
% INFO: Zeilenabstand setzen:
%
% Befehle:
% - \singlespacing  => 1-zeilig (Standard)
% - \onehalfspacing => 1,5-zeilig
% - \doublespacing  => 2-zeilig 
\onehalfspacing % Zeilenabstand auf 1,5-zeilig setzen

% Farbboxen (für die Merkkästen in dieser Vorlage):
\usepackage{tcolorbox}
\tcbset{colback=white,colframe=orange,
        fonttitle=\bfseries}

\usepackage[colorlinks,pdfpagelabels,pdfstartview=FitH,
bookmarksopen=true,bookmarksnumbered=true,linkcolor=black,
plainpages=false,hypertexnames=false,citecolor=black]{hyperref} % Für Verlinkungen
% INFO: Verlinkungen mit dem hyperref-Paket:
%
% Die Angabe von URLs mit dem Befehl \url{} erlaubt einen
% gesonderten Umgang mit Weblinks. Denn die Links werden verlinkt.
% Auch erfolgt automatisch am Zeilenende ein Umbruch des Links.
% Es ist auch nicht erforderlich, Sonderzeichen in der URL manuell zu 
% entschärfen.
%
% TIPP: Sollte ein Umbuch bei einem Link nicht automatisch erfolgen, so kann
% das daran liegen, dass ein/mehrere Zeichen zusätzlich angegeben werden müssen,
% an dem der Link umbrochen werden kann.
% Dies kann mit folgendem Befehl erfolgen (Beispiel):
% \renewcommand*\UrlBreaks{\do-\do_}

% Das Paket "biblatex" für autom. 
% Literaturverzeichnisse:
\usepackage{csquotes} % Für sprachangepasste Anführungszeichen
\usepackage[backend=bibtex,style=alphabetic]{biblatex}
\addbibresource{bib/literatur.bib}           

%%%%%%%%%%%%%%%%%%%%%%%%%%%%%%%%%%%%%%%%%%%%%
%% DOKUMENT                                %%
%%%%%%%%%%%%%%%%%%%%%%%%%%%%%%%%%%%%%%%%%%%%%
\begin{document}
  
  % Deckblatt
  \pagestyle{empty}
  \begin{titlepage}
    \includegraphics[scale=1.00]{Sources/logo_TH-Koeln_CMYK_22pt}\\
    \begin{center}
      \Large
      Technische Hochschule Köln\\
      Fakultät für Informatik und Ingenieurwissenschaften\\
      \hrule\par\rule{0pt}{2cm} % Horizontale Trennlinie  mit 2 cm Abtand nach unten erzeugen
      \LARGE
      \textsc{Bachelorarbeit - Expose}\\
      \vspace{1cm} % Vertikaler Abstand von 1cm erzeugen
      \huge
      Der Einfluss von Farbnormalisierung in neuronalen Netzen\\
      \vspace{1 cm}
      \large
      Vorgelegt an der TH Köln\\
      Campus Gummersbach\\
      im Studiengang\\
      Medieninformatik\\ 
      \vspace{1.0cm}
      ausgearbeitet von:\\
      \textsc{Torben Krause}\\
      (Matrikelnummer: 11106885)\\
      \vspace{1.5cm}
      \begin{tabular}{ll} % Einfache Tabelle ohne Rahmen, mit 2 Spalten erzeugen
          \textbf{Erster Prüfer:} & Prof. Dr. Martin Eisemann \\
          \textbf{Zweiter Prüfer:} & Prof. Dr. ????? \\
      \end{tabular}
      \vspace{1.5cm}
      \\Gummersbach, im Februar
    \end{center}    
  \end{titlepage}
    
  \newpage
  \pagestyle{plain}
  \pagenumbering{arabic}
  \setcounter{page}{1}
  
  \tableofcontents
  
  \newpage
  
  \section{Problemstellung}
  Das Trainieren eines künstlichen neuronalen Netzes ist ein langwieriger Prozess, bei dem 
  eine große Menge von Ressourcen benötigt werden. Zum einen werden Hardwarekomponenten benötigt, 
  welche für Matrizenmultiplikation geeignet sind, und zum anderen wird je nach Umfang des 
  neuronalen Netzes eine gewisse Trainingszeit vorausgesetzt. Zusätzlich werden für ein 
  effektives neuronales Netz oftmals mehrere tausend Trainingsdaten benötigt, welche zusätzlich 
  durch ihre Größe die Trainingszeit beeinflussen. Neben der Größe der Bilder spielt die  
  Ausleuchtung der Objekte auf den Bildern eine große Rolle. Durch eine einseitige Beleuchtung 
  entstehen Schatten und Erhellungen, welche die Farben der Objekte verändern. Durch 
  die verfälschten Farbwerte können Probleme bei der richtigen Zuordnung der Objekte entstehen. 
  Diese können die richtige Erkennungsrate
des künstlichen neuronalen Netzes beeinflussen, da mehrere Farbwerte für die selbe Oberfläche möglich sind. Das kann zu fehlerhaften Zuweisungen zwischen den erlernten Klassen führen.
  \section{Lösungsansatz}
Eine Möglichkeit um die Trainingsergebnisse optimieren zu können besteht darin, die Farbwerte der Trainingsdaten zu normalisieren, um mögliche Farbverfälschungen, welche durch unausgeglichene Belichtungen entstehen, auszugleichen. Ursachen für Farbverfälschungen sind beispielsweise, unterschiedliche Licht-Farben/Positionen oder auch mehrere Lichtquellen. Bei einer Farbnormalisierung, werden beispielsweise benachbarte Pixel, die im gleichen Farbraum liegen, auf einen Farbwert angepasst, damit keine großen Farbabweichungen auftreten. Es gibt viele derartige Algorithmen, welche unterschiedliche Ansätze verfolgen. Um einen Einblick in die Möglichkeiten der Algorithmen zu zeigen, werden einige im folgenden beschrieben:
  \begin{itemize}
  \item Die umfassende Farbnormalisierung versucht die Lokalisierungs- und Objektklassifizierungsergebnisse, in Kombination mit der Farbindizierung, zu erhöhen \cite{finlayson1998comprehensive}. Dieser Algorithmus ist iterativ und arbeitet in zwei Stufen. Die erste Stufe besteht darin, den roten, grünen und blauen Farbraum mit normalisierter Intensität zu verwenden, um jedes Pixel zu normalisieren. Die zweite Stufe sorgt für eine Normalisierung jedes Farbkanals separat, wodurch die Summe der Farbkomponenten einem Drittel der Pixelanzahl entspricht. Durch diese Normalisierung wird die Anzahl der Farbwerte in einem Bild verringert.
  \item Ein weiterer Ansatz ist der Gray World Algorithmus \cite{buenaposada2001variations}. Bei dieser Normalisierung, wird davon ausgegangen, das es drei Faktoren gibt, welche das Beleuchtungsspektrum verändern können. Diese werden auf den roten, grünen und blauen Farbkanälen angewendet. Eine Veränderung der beleuchteten Farbe kann durch die Skalierung von $\alpha$, $\beta$ und $\gamma$ in den RGB-Farbkarnälen modelliert werden. Es kann eine Konstanzlösung erzielt werden, indem jeder Farbkanal durch seinen Durchschnittswert geteilt wird.
  \item Die Histogram Equalization ist eine nicht lineare Transformation, welche die Pixelstruktur beibehält und in der Lage ist, eine monoton steigende Farbtransformation zu Normalisieren \cite{goatman2003colour}. Es gilt als eine stärkere Methode als der Gray-World-Algorithmus. Die Ergebnisse dieser Methode, haben üblicher weise einen dominanten blauen Farbkanal, welcher das Bild unnatürlich aussehen lässt.
  \end{itemize}
Die aufgeführten Verfahren stellen nur einen Teil der Algorithmen, welche verwendet werden können dar. Für die Bachelorarbeit sollen im vorhinein einige der Verfahren ausgewählt und getestet werden.
  \section{Vorgehen}
Für die Bachelorarbeit sollen verschiedene Bildnormalisierungs-Algorithmen auf verschiedenen Datensätzen angewendet und im Training verwendet werden. Diese Datensätze bestehen aus mehreren Klassen welche mit unterschiedlichen Lichtverhältnissen aufgenommen werden. Daraufhin werden die verschiedenen Netze miteinander verglichen. Dabei soll herausgearbeitet werden, welchen Einfluss die normalisierten Bilder auf das Training haben. Bestenfalls soll die These bestätigt und eine Erhöhung der Effizienz des Trainings erkannt werden. Weiterhin soll geschaut werden, welchen Einfluss eine Verringerung der Trainingsdaten auf das neuronale Netz haben. Sollte eine Normalisierung von Daten die Effizienz des Trainings erhöhen, könnten Ressourcen eingespart werden.\\\\
\textbf{Geplante Durchführung}\\\\
(Die Schritte 2-5 werden für jedes ausgewählte Verfahren durchgeführt. Jedes Training wird auf unterschiedlichen Netzen durchgeführt.)
  \begin{enumerate}
\item Training der Datensätze ohne Normalisierung
\item Datensatz Farbnormalisieren
\item Training der Normalisierten Datensätzen auf einem weiteren Netz
\item Auswertung des Trainings
\item Vergleich des Normalisierungs-Verfahren
  \end{enumerate}
  
  \newpage
  
  \section{Vorläufige Gliederung}
  
  \begin{enumerate} %Gliederung
  \item Deckblatt
  \item Vorwort
  \item Danksagung
  \item Abstract
  \item Inhaltsverzeichnis
  \item Abbildungs- und Tabellenverzeichnis
  \item Abkürzungsverzeichnis
  
  \item Einleitung
  \begin{enumerate} %Einleitung
  \item Problemstellung und Ziele
  \item Begründung der These
  \item Anforderungen
  \item Struktur der Arbeit
  \end{enumerate} %Einleitung
  \item Theoretische Grundlagen
  \begin{enumerate} %Theoretische Grundlagen
  \item neuronale Netze
  \begin{enumerate} %neuronale Netze
  \item Funktionsweise künstlicher neuronaler Netze
  \item Entstehung von Trainingsdaten
  \end{enumerate} %neuronale Netze
  \item Das digitale Bild
  \begin{enumerate} %Das digitale Bild
  \item Aufbau vom Bildern
  \item Einfluss von Belichtungen
  \item Bild-Normalisierung
  \begin{enumerate}
  \item Umfassende Farbnormalisierung 
  \item Gray-World-Algorithmus
  \item Histogram Equalisation
  \end{enumerate}
  \end{enumerate} %Das digitale Bild
  \end{enumerate} %Theoretische Grundlagen
  \item Methodik und Durchführung
  \item Auswertung der Ergebnisse
  \item Diskussion
  \item Fazit
  
  \item Literaturverzeichnis
  \item Anhang
  \item Eidesstattliche Erklärung
  \end{enumerate} \vspace{2 cm} %Gliederung
 
  
  \section{Zeitplan}  
  
  \begin{enumerate}
  \item Ab 11.02:     	Vorbereitung und Recherche
  \item Start ??.??: 	Bachelorarbeit
  \item 7 Tage: 		Literaturrecherche
  \item 21 Tage: 		Durchführung des praktischen Teils + Auswertung der Trainingsdaten
  \item 14 Tage:		Rohfassung Hauptteil
  \item 7 Tage:		Rohfassung Einleitung + Schluss
  \item 14 Tage:		Überarbeitung + Korrektur
  \item 4 Tage:		Layout + Titelblatt
  \item 3 Tage:		Druck + Binden
  \item 2 Tage:		Abgabe
  \end{enumerate}
  
  \newpage
  
  \printbibliography
  
\end{document}

